\chapter{串联电路}\label{ch:g4:series-circuits}

去年，一串彩灯因为缺了一只
灯泡而全部熄灭——我们当时承诺过，这种连接方式有个名字。
名字就是串联电路：把许多用电器接在同一个回路上，这是最简单的办法。
它有严格的规则，你抽屉里的手电筒
一直都在遵守它们。

\section{一个回路，多个用电器}

\begin{definition}[串联]\label{def:g4:series-circuits:series}
电路中的用电器若\emph{串联}\index{串联电路}，
即它们一个接一个地排列在同一个回路上，
就像一串项链上的珠子：电流的路径依次经过第一个、第二个、第三个，
中间没有任何分岔。按这种方式连接的电路，
称为\emph{串联电路}。
\end{definition}

\begin{omfigure}
\begin{tikzpicture}[scale=0.9]
  \draw (0,0) to[battery1, l=battery] (0,2.6)
    to[short] (1.6,2.6)
    to[cspst, l=switch] (3.2,2.6)
    to[lamp, l=bulb 1] (5.2,2.6)
    to[lamp, l=bulb 2] (7.2,2.6)
    to[short] (7.2,0)
    to[short] (0,0);
\end{tikzpicture}

{\small 串联电路：电池、开关和两只灯泡位于同一个回路上——
每个用电器都穿在同一根项链上。}
\end{omfigure}

\begin{method}[搭建电路]\label{met:g4:series-circuits:build}
用一节电池、两只灯泡、一个开关和若干导线，按以下步骤连接：
\begin{enumerate}
  \item 用导线把电池的 $+$ 接线柱连到
        开关；
  \item 用导线从开关连到第一只灯泡，再从第一只灯泡
        连到第二只灯泡；
  \item 闭合回路：从第二只灯泡连回电池的
        $-$ 接线柱；
  \item 合上开关，观察两只灯泡是否同时亮起。
\end{enumerate}
开灯前，先用手指沿回路走一遍：如果手指
需要在两条路之间做选择，那就不是串联电路。
\end{method}

\section{串联电路的规则}

\begin{proposition}[单一回路的规则]\label{prop:g4:series-circuits:rules}
在串联电路中：
\begin{enumerate}
  \item \emph{全有或全无}：任何一处断开——开关断开，
        灯泡损坏，导线松脱，所有用电器便
        同时停止工作；
  \item \emph{顺序无关}：交换两只灯泡，或把
        开关移到回路的另一侧，一切都与原来
        完全一样；
  \item \emph{分担}：回路上的灯泡越多，
        每只灯泡越暗；电池的“推力”
        被整串灯泡分担。
\end{enumerate}
\end{proposition}

\begin{example}[观察规则]\label{ex:g4:series-circuits:seeing}
搭建双灯泡电路，逐条检验规则。拧下\emph{任一只}
灯泡：两只灯泡都熄灭——这是规则一（被拧下的灯座相当于断点）。
交换两只灯泡，或把开关移到二者之间：
亮度毫无变化——这是规则二。现在改用一只灯泡重新搭建，它
独自发光时比有“伴侣”时明显更亮——
这正是规则三的反向体现。
\end{example}

\begin{example}[彩灯之谜的解答]\label{ex:g4:series-circuits:partylights}
旧日的谜团如今已不再是谜：彩灯上的灯泡是串联的，
几十颗珠子串在同一根项链上。缺一只灯泡就是一处断点，
规则一使整串彩灯全部熄灭。逐只试灯泡找出坏的那只，
曾让几代家庭用最笨的办法领会规则一——
现代彩灯串的结构已经不同，
下一年会讲到原因。
\end{example}

\begin{remark}[开关控制所有用电器]\label{rem:g4:series-circuits:switch}
规则二有一个方便的推论：一个开关，无论放在回路
的哪个位置，都能控制回路上的所有用电器。这就是为什么一个墙上的
开关可以控制整排的天花板灯——
同样的道理，你无法只关掉串联灯串中的\emph{一只}灯泡，
否则会让全部灯泡一起熄灭。
\end{remark}

\section{串联更多电池}

\begin{example}[首尾相接]\label{ex:g4:series-circuits:batteries}
电池也可以串联排列——不过它们比灯泡更挑剔：
它们必须\emph{首尾相接}：每节电池的 $+$
接在下一节电池的 $-$ 上。两节电池这样排列时，推力
相加，灯泡明显比用一节时更亮。把其中一节
电池调个头，变成头对头——两股推力
相互抵消，灯泡不亮。
\end{example}

\begin{omfigure}
\begin{tikzpicture}[scale=0.8]
  \draw (0,0) to[battery1] (0,1.8)
    to[battery1, l=two batteries] (0,3.6)
    to[short] (3.6,3.6)
    to[lamp, l=brighter!] (3.6,0)
    to[short] (0,0);
  \begin{scope}[xshift=6.9cm]
    \draw (0,0) to[battery1, l=one battery] (0,1.8)
      to[short] (0,3.6)
      to[short] (3.6,3.6)
      to[lamp, l=dimmer] (3.6,0)
      to[short] (0,0);
  \end{scope}
\end{tikzpicture}

{\small 两节首尾相接排成一列的电池共同推动：
同一个灯泡比只用一节电池时更亮。}
\end{omfigure}

\begin{example}[手电筒内部]\label{ex:g4:series-circuits:flashlight}
打开手电筒，电池一节接一节地滑入
筒中——正好是一列首尾相接的队伍——然后灯泡和
按钮开关沿外壳连成一个单一回路。
手电筒\emph{就是}口袋里的串联电路。现在你明白了
为什么只要有一节电池装反，手电筒就不亮：
电池头对头，推力相互抵消。
\end{example}

\section{练习}

\begin{exercise}[$\star$]\label{exo:g4:series-circuits:1}
“串联”是什么意思？本章借用日常生活中的哪个
情景来说明它？
\end{exercise}

\begin{exercise}[$\star$]\label{exo:g4:series-circuits:2}
说出串联电路的三条规则。
\end{exercise}

\begin{exercise}[$\star$]\label{exo:g4:series-circuits:3}
在双灯泡电路中，灯泡 2 被拧下，灯泡 1 会怎样，
为什么？这里起作用的是哪条规则？
\end{exercise}

\begin{exercise}[$\star$]\label{exo:g4:series-circuits:4}
佐伊把开关从灯泡 1 前面移到灯泡 2 后面。
电路的工作情况会发生什么变化？依据哪条规则？
\end{exercise}

\begin{exercise}[$\star$]\label{exo:g4:series-circuits:5}
一节电池能让一个灯泡正常发光。预测下列情况下的亮度：
两个灯泡串联；\ 三个灯泡串联。被共享的是什么？
\end{exercise}

\begin{exercise}[$\star$]\label{exo:g4:series-circuits:6}
两节电池应怎样放置才能共同推动？画出电池队列，
并在每节电池上标出 $+$ 和 $-$。
\end{exercise}

\begin{exercise}[$\star$]\label{exo:g4:series-circuits:7}
沿手电筒的回路追踪：说出电流经过的各个元件，
从电池的 $+$ 极出发，回到它的 $-$ 极。
\end{exercise}

\begin{exercise}[$\star$]\label{exo:g4:series-circuits:8}
门铃电路中，电池、按钮和铃串联在一起。
在哪里可以再加一个按钮，使铃只有在
\emph{两个}按钮同时按下时才响？
\end{exercise}

\begin{exercise}[$\star$]\label{exo:g4:series-circuits:9}
一串十只串联的彩灯不亮了。灯泡看上去都
完好。请描述一个有耐心的方案，用一只确定完好的灯泡
去找出坏掉的那一只。
\end{exercise}

\begin{exercise}[$\star\star$]\label{exo:g4:series-circuits:10}
汤姆的手电筒用两节电池，发光很亮。在
黑暗中更换电池后，手电筒却不亮了——可两节新电池都是满电的，
灯泡也完好，导线也没有松动。汤姆做了什么，
手电筒为什么仍然不亮？（见
\cref{ex:g4:series-circuits:batteries}）
\end{exercise}

\begin{exercise}[$\star\star$]\label{exo:g4:series-circuits:11}
一位工程师必须为一台机器装一个紧急停止开关，
这台机器的电动机与其电池接在同一个单回路上。她说：
“回路上任何位置都可以。”她说得对吗？
如果这台机器的电路有支路，她的方案为什么会失败？（下
学年的章节会给它们起一个名字。）
\end{exercise}
